
\section{Auxiliary uniformity details}
\label{app:uniformity}

This appendix records several deterministic and probabilistic steps used in
the main text.

\subsection{Chronological likelihood cancellation}

For a fixed horizon, disintegrate the joint law in the order
\[
 A_1,Y_1,A_2,Y_2,\ldots,A_n,Y_n.
\]
Under a parameter-independent policy rule, the conditional density of \(A_i\)
given the realized past is the same function under all parameter values.
Consequently its factor cancels in a likelihood ratio at that record.  The
observation factors do not cancel.  The resulting candidate state and
candidate mean must still be recomputed at each parameter while holding the
realized action record fixed.  No conditioning on the entire future action
sequence is used.  If the policy rule itself depended on the candidate
parameter, the conditional action Radon--Nikodym factors would have to remain
in the likelihood; such policies are excluded from Definitions
\ref{def:triangular} and the two mechanical policy classes.

\subsection{Derivative formulas for normalized measures}

Let \(V_\theta\) be a bounded \(C^r\) function on a measurable space and
\[
 w_\theta(dx)=
 \frac{e^{V_\theta(x)}\nu(dx)}{\int e^{V_\theta}d\nu}.
\]
Then, as signed measures,
\begin{align*}
 \partial_jw_\theta
 &=\{\partial_jV_\theta-
       w_\theta(\partial_jV_\theta)\}w_\theta,\\
 \partial_{jk}w_\theta
 &=\bigl[(\partial_jV-\bar V_j)(\partial_kV-\bar V_k)
   +\partial_{jk}V-w(\partial_{jk}V)
   -w\{(\partial_jV-\bar V_j)(\partial_kV-\bar V_k)\}\bigr]w_\theta,
\end{align*}
where \(\bar V_j=w_\theta(\partial_jV)\).  Induction shows that every
higher derivative is a polynomial in centered derivatives of \(V\),
multiplied by \(w_\theta\).  Thus
\[
 \norm{\partial^\alpha w_\theta}_{TV}
 \le C_r\left(1+\max_{\abs\gamma\le r}
 \norm{\partial^\gamma V_\theta}_\infty\right)^{C_r},
 \qquad \abs\alpha\le r.
\]
The same formulas, together with uniform Banach-space Taylor remainders for
\(V_{\theta+h}\), prove strong total-variation differentiability.  No density
of \(\nu\) is needed.

\subsection{Interpolation constants for the six-duration map}

With the notation of Section~\ref{sec:lattice}, the exact interpolation
identity is
\[
 (y_1,\ldots,y_6)^T
 =VD_\tau(r_2,\ldots,r_7)^T+\mathcal E_\tau.
\]
The \(j\)-th coordinate of \(\mathcal E_\tau\), together with every first
parameter derivative, is bounded by \(C(j\tau)^9\).  Hence
\[
 \norm{D_\tau^{-1}V^{-1}\mathcal E_\tau}
 \le \norm{D_\tau^{-1}}\norm{V^{-1}}C(6\tau)^9
 \le C_V\tau.
\]
This is the claimed \(C^1\) error.  The determinant factorization is
\[
 \det V=\left(\prod_{j=1}^6j^3\right)
 \left(\prod_{r=3}^8\frac1{r!}\right)
 \prod_{1\le i<j\le6}(j-i),
\]
which is nonzero.

\subsection{The compact full-index net}

The radial set
\[
 \mathscr K=\{(\theta_0,r,v):\theta_0\in\Theta,\ v\in S^{p-1},
 r\ge0,\ \theta_0+rv\in\Theta\}
\]
is compact because it is closed and \(r\) is bounded by the diameter of
\(\Theta\).  For
\[
 q_i(\theta_0,r,v)=\int_0^1v^TDm_i^{\theta_0+srv}ds,
\]
the candidate response derivatives make \(q_i\) continuous at \(r=0\) and
uniformly Lipschitz in all three variables.  The square has Lipschitz constant
at most twice the product of the value and first-index derivative bounds.
The same deterministic constant applies to
\(\E(q_i^2\mid\cF_{i-1})\), since conditional expectation is a contraction in
supremum norm.  A single deterministic finite net therefore controls both the
realized and predictable fields.  Its cardinality depends on the parameter
box and mesh, but not on \(n\), the policy, the true parameter, or the true
state.

\subsection{Relative normalization in total variation}

Suppose \(f_n,g_n\ge0\), \(G_n=\int g_n>0\), and
\[
 \frac{\int\abs{f_n-g_n}}{G_n}\to0.
\]
Then \(\abs{F_n-G_n}/G_n\to0\), where \(F_n=\int f_n\), and eventually
\(F_n/G_n\) is bounded away from zero.  Therefore
\begin{align*}
 \int\left|\frac{f_n}{F_n}-\frac{g_n}{G_n}\right|
 &\le \frac1{F_n}\int\abs{f_n-g_n}
   +\frac{\abs{F_n-G_n}}{F_n}\to0.
\end{align*}
Half of this integral is the total-variation distance.  This elementary step
is why Lemma~\ref{lem:random-laplace} proves a relative unnormalized
\(L^1\) estimate rather than only pointwise convergence of log likelihoods.

\subsection{Uniform strong law for the countable-duration experiment}

Let \(\mathcal D\) be the compact image of \(\mathfrak B\) in
\(\ell^2(\omega)\), with uniformly bounded coordinates.  For any
\(\varepsilon>0\), choose a finite \(L^2(\omega)\) net
\(d^1,\ldots,d^N\).  At a fixed net point,
\[
 M_n(d)=\sum_{i=1}^nS_i\xi_id(M_i)
\]
is a square-integrable martingale with
\(\E M_n(d)^2\le Cn\norm d_{L^2(\omega)}^2\).  The ordinary martingale strong
law gives \(M_n(d)/n\to0\).  For an approximation error \(e=d-d^k\), Doob's
inequality gives
\[
 \Prob\left\{\max_{m\le2^r}\abs{M_m(e)}>2^r\delta\right\}
 \le C\varepsilon^2/(2^r\delta^2).
\]
The sum over \(r\) is finite.  Borel--Cantelli and the finite net show
\[
 \sup_{d\in\mathcal D}\abs{M_n(d)}/n\to0
\]
almost surely after \(\varepsilon\downarrow0\) along a sequence.  Apply the
same argument to the uniformly bounded square class, using
\[
 \abs{x^2-y^2}\le2B\abs{x-y},
\]
to obtain the uniform square strong law in
\eqref{eq:uniform-square-slln}.

\subsection{Analytic continuation from a finite duration interval}

For every bounded generator \(A_\beta\), the map
\[
 t\mapsto h_\beta(t)=\ell\int_0^te^{sA_\beta}Bds
\]
is entire.  If two such responses agree on the dense diagnostic set in
\((0,T]\), continuity makes them equal on that interval, and the identity
theorem makes them equal for every real \(t\ge0\).  This justifies the use of
the complete Weyl reconstruction after identification on a finite duration
interval.
